\section{Hamiltonian Ising sebagai Fondasi Masalah Pasar Finansial}

\subsection{Representasi Kualitatif dalam Pasar Finansial}
Dalam domain ekonofisika, agen ekonomi secara formal dianalogikan sebagai sistem partikel diskrit yang saling berinteraksi secara kolektif di bawah pengaruh parameter eksternal pasar yang dinamis. Setiap agen direpresentasikan sebagai entitas \textit{spin} ($s_i$) yang memiliki dua arah keadaan tetap sebagai manifestasi dari keputusan biner dalam pengambilan posisi investasi. Keadaan $+1$ merepresentasikan posisi beli (\textit{long}) yang mencerminkan optimisme pasar, sementara keadaan $-1$ merepresentasikan posisi jual (\textit{short}) yang menandakan adanya tekanan penawaran. Figure 7 mengilustrasikan pemetaan teknis antara arah orientasi \textit{spin up/down} dengan posisi strategis \textit{buy/sell} di dalam pasar finansial untuk mempermudah visualisasi dinamika keputusan biner. Representasi fisik ini memungkinkan dilakukannya analisis mendalam terhadap munculnya tren makroekonomi yang persisten berdasarkan agregasi keputusan mikro yang diambil oleh setiap individu agen secara stokastik.

Fenomena perilaku kawanan (\textit{herding behavior}) di pasar modal terjadi secara sistemik saat interaksi antar agen menjadi sangat dominan sehingga memicu pergerakan harga yang searah secara masif. Secara fisik, kondisi ini setara dengan fase teratur atau \textit{ferromagnetic coupling} di mana mayoritas \textit{spin} dalam sistem menunjuk ke arah yang sama akibat kuatnya gaya tarik antar agen. Sebaliknya, aliran berita fundamental atau perubahan kebijakan moneter eksternal berperan sebagai medan magnet luar (\textit{external field}) yang memberikan bias probabilitas pada arah pilihan strategi setiap agen. \textit{Hamiltonian} sistem kemudian diimplementasikan untuk menghitung total konflik atau keselarasan energi dalam konfigurasi pasar tersebut guna menentukan stabilitas strategi kolektif yang diadopsi oleh para pelaku pasar \cite{galam_ising_2010}.

\subsection{Pengenalan Hamiltonian Ising}
Model \textit{Ising} awalnya dikembangkan di dalam bidang mekanika statistik guna mendeskripsikan fenomena transisi fase pada material magnetik melalui interaksi mikroskopis antar derajat kebebasan yang bersifat diskrit. Model ini mengandalkan fungsional energi atau \textit{Hamiltonian} sebagai instrumen utama untuk menentukan tingkat stabilitas dari setiap kemungkinan konfigurasi partikel di dalam sebuah sistem kisi. Figure 8 menyajikan representasi struktur \textit{grid} dari model \textit{Ising} yang memperlihatkan interaksi kopling $J$ antar tetangga terdekat serta pengaruh medan magnet eksternal $h$ terhadap orientasi setiap partikel. Fleksibilitas struktur matematika ini menjadikannya sangat kuat untuk memetakan berbagai masalah optimasi kombinasional yang kompleks ke dalam sebuah lanskap energi fisik yang dapat diselesaikan secara efisien.

Persamaan \textit{Hamiltonian} secara teknis menghitung total biaya energi sistem berdasarkan keselarasan orientasi antar pasangan \textit{spin} yang berinteraksi di dalam ruang parameter yang ditentukan. Apabila parameter interaksi feromagnetik memiliki nilai $J > 0$, maka sistem akan memberikan nilai energi yang lebih rendah dan lebih stabil saat para agen memilih posisi yang sama. Medan magnet eksternal $h$ akan memberikan gaya dorong pada setiap \textit{spin} individu agar selaras dengan arah tren dominan atau kebijakan otoritas yang berlaku di pasar. Melalui mekanisme operasional ini, proses pencarian keputusan investasi yang optimal secara matematis menjadi setara dengan pencarian konfigurasi \textit{ground state} atau energi minimum di dalam sistem fisik tersebut.
\begin{equation}
    H(\sigma) = -\sum_{i < j} J_{ij} \sigma_i \sigma_j - \sum_{i=1}^{N} h_i \sigma_i
\end{equation} (1)

\subsection{Formulasi Objektif Markowitz pada Portofolio Biner}
Strategi diversifikasi fundamental yang dirumuskan oleh Markowitz bertujuan secara spesifik untuk mencari titik keseimbangan optimal antara ekspektasi imbal hasil maksimal dengan risiko volatilitas yang seminimal mungkin. Dalam skenario portofolio biner, investor dihadapkan pada keputusan mutlak untuk menentukan apakah suatu aset tertentu akan dimasukkan ke dalam keranjang investasi atau dikeluarkan sepenuhnya. Pergeseran variabel dari domain kontinu ke domain biner ini mengubah masalah pemilihan portofolio menjadi persoalan optimasi kombinasional yang secara komputasi termasuk ke dalam kelas \textit{NP-hard}. Tantangan teknis utama dalam formulasi ini adalah besarnya ruang solusi yang tumbuh secara eksponensial seiring dengan bertambahnya jumlah aset $N$, yang mencapai total $2^N$ kemungkinan kombinasi strategis.

Fungsi biaya Markowitz secara sistemik menggabungkan elemen matriks kovarians dan ekspektasi imbal hasil melalui parameter penghindaran risiko (\textit{risk aversion}) guna menghasilkan metrik efisiensi portofolio tunggal. Matriks kovarians ($\Sigma_{ij}$) di dalam rumus ini merepresentasikan interaksi risiko sistemik antar pasangan aset yang menentukan derajat diversifikasi di dalam satu keranjang investasi. Struktur fungsi objektif ini bersifat kuadratik, yang secara matematis memiliki bentuk isomorfisma yang identik dengan interaksi berpasangan antar partikel magnetik di dalam model \textit{Ising} \cite{lucas_ising_2014}. Kesamaan struktural ini memungkinkan masalah seleksi aset yang sangat kompleks untuk diselesaikan secara langsung melalui mekanisme simulasi sistem fisik atau algoritma optimasi berbasis energi.
\begin{equation}
    f(x) = \lambda \sum_{i=1}^{N} \sum_{j=1}^{N} \Sigma_{ij} x_i x_j - \sum_{i=1}^{N} \mu_i x_i
\end{equation} (2)

\subsection{Model Mean-Variance Portfolio Optimization}
Model \textit{Mean-Variance} bekerja secara operasional dengan mengidentifikasi titik-titik \textit{Pareto Optimum} pada ruang dua dimensi yang dibentuk oleh parameter risiko dan imbal hasil aset. Titik-titik optimal ini membentuk sebuah kurva yang dikenal sebagai \textit{Efficient Frontier}, yang menjadi acuan standar bagi setiap investor rasional dalam melakukan alokasi kapital yang efisien. Figure 9 mengilustrasikan plot kurva \textit{Efficient Frontier} beserta posisi titik \textit{Pareto Optimum} yang mewakili kombinasi portofolio terbaik untuk setiap tingkat toleransi risiko tertentu. Melalui kurva ini, investor dapat dengan mudah memilih kombinasi aset yang menawarkan risiko terendah untuk target keuntungan tertentu atau keuntungan tertinggi pada batas risiko tertentu.

Isomorfisme antara model \textit{Mean-Variance} dan sistem \textit{Magnetic Ising} terletak pada kesetaraan fungsional antara parameter ekonomi dengan variabel fisik magnetik yang mendasarinya. Suku varians portofolio yang merepresentasikan risiko sistemik memiliki peran fungsional yang sama dengan energi interaksi antar \textit{spin} ($J$) di dalam sistem partikel magnetik. Sementara itu, variabel ekspektasi imbal hasil aset berperan sebagai medan magnet luar ($h$) yang memberikan gaya penggerak terhadap orientasi setiap \textit{spin} individu di dalam sistem. Hubungan isomorfisma ini membuktikan bahwa mekanisme pasar yang efisien secara intrinsik mengikuti prinsip alamiah dalam meminimalkan energi total sistem guna mencapai stabilitas ekuilibrium \cite{markowitz_portfolio_1952}.

\subsection{Pemetaan Hamiltonian Ising}
Untuk mengeksekusi proses optimasi pada perangkat keras komputer kuantum, variabel keputusan biner klasik harus dipromosikan menjadi operator mekanika kuantum yang bekerja di dalam ruang Hilbert. Variabel keputusan biner dalam masalah portofolio kemudian dipetakan secara spesifik ke dalam operator \textit{Pauli-Z} ($\hat{Z}$) untuk merepresentasikan status keputusan pasar secara digital. Nilai eigen dari operator ini, yaitu $+1$ dan $-1$, secara alami mewakili status basis \textit{qubit} $|0\rangle$ dan $|1\rangle$ yang merepresentasikan keputusan beli atau jual. Pemetaan ini menjamin bahwa interaksi kovarians antar aset dapat diterjemahkan menjadi korelasi fase antar \textit{qubit} yang saling terikat secara kuantum di dalam sirkuit optimasi.

Penurunan parameter fisik akhir dilakukan melalui substitusi matematis menggunakan transformasi variabel $x = (1 - \sigma)/2$ guna menjembatani format QUBO dengan format \textit{Ising}. Substitusi ini menghasilkan rumusan eksplisit bagi medan lokal ($h_i$) dan kekuatan interaksi ($J_{ij}$) yang akan dikonfigurasi pada perangkat keras kuantum sebagai parameter Hamiltonian. Dengan parameter hasil pemetaan ini, algoritma seperti \textit{Variational Quantum Eigensolver} (VQE) dapat mulai mengeksplorasi ruang solusi melalui manipulasi fase dan keterikatan antar \textit{qubit} secara simultan. Keberhasilan transformasi ini memberikan jaminan bahwa setiap hasil pengukuran pada \textit{qubit} memiliki interpretasi ekonomi yang sah dan dapat diimplementasikan kembali sebagai strategi alokasi portofolio riil.
\begin{equation}
    \hat{H} = \sum_{i=1}^{N} h_i \hat{Z}_i + \sum_{i < j} J_{ij} \hat{Z}_i \hat{Z}_j
\end{equation} (3)
